\specialchapter{术语索引}

参考编号依次表示章节、节和小节（或习题）。

阿代尔、限制阿代尔、主限制阿代尔：VII.2.4

m-进滤过：III.2.1

--- 拓扑：III.2.5

\(n\)-进整数：III.2.12

Azumaya 代数：II.1.7

--- 有限生成的：III.1.1

--- 在环上整且有限的：V.1.1

域在代数中的代数闭包：V.1.2

在代数中代数闭的域：V.1.2

--- 代数相关、无关的元素：III.1.1

--- 代数自由、相关的族：III.1.1

对几乎所有 \(\mathfrak{p} \in \mathbf{P}(\mathbf{A})\)（成立的性质）：VII.4.3

--- 殆幂零自同态：IV.1.4

绝对值逼近定理：VI.7.3

--- 赋值逼近定理：VI.7.2

Artin-Rees 引理：III.3.1

与分次模相伴的（滤过模）：III.2.1

--- 与分次环相伴的（滤过环）：III.2.1

--- 与分次相伴的（滤过）：III.2.1

--- 与滤过相容的同态相伴的分次同态：III.2.4

--- 与滤过模相伴的（分次模）：III.2.3

--- 与滤过环相伴的（分次环）：III.2.3

--- 与环同态相伴的（映射）：II.4.3

--- 与模相伴的（素理想）：IV.1.1

--- 伴随环、位、赋值：VI.3.3

位的典范分解：VI.2.3

--- 赋值的典范分解：VI.3.2

--- 从素理想 \(\mathfrak{p}'\) 的分解群到

A'/p' 的自同构群的典范同态：V.2.2

与有限生成模相联的除子类：VII.4.7

除子类（除子幺半群）：VII.1.2

域在代数中的代数闭包：V.1.2

--- 整环的整闭包：V.1.2

--- 环在代数中的整闭包：V.1.2

交换图：I.1.2

（滤过）与环结构、模结构相容：III.2.1

--- （同态）与滤过相容：III.2.4

赋值扩张的完备系：VI.8.2

完全整闭整环：V.1.4

不可约分支（拓扑空间的）：II.4.1

Hensel 条件：III.4.5

子模的导子：V.1.5

伪 Bézout 整环上多项式的容度：VII.1.习题 23

--- 扭模的容度：VII.4.5

Eisenstein 不可约判别法：VII.3.习题 20

位的典范分解：VI.2.3

--- 素理想的完全分解：V.2.2

--- 素理想的分解域：V.2.2

--- 素理想的分解群、分解环：V.2.2

--- Dedekind 整环中理想分解为素因子：VII.2.3

--- 准素：IV.2.2 和习题 20

--- 既约准素：IV.2.3 和习题 20

递减滤过：III.2.1

Dedekind 整环：VII.2.1

（拓扑）由滤过定义：III.2.5

定义理想：III.3.2

一个赋值相对于另一个赋值的剩余类次数：VI.8.1

整依赖（方程）：V.1.1

（模滤过）由环滤过导出：III.2.1

交换图：I.1.2

--- 蛇形图：I.1.4

离散滤过：III.2.1

--- 离散赋值：VI.3.6

特异多项式：VII.3.8

除子、主除子：VII.1.1

--- 行列式除子：VII.4.习题 11

--- 有限生成除子：VII.1.习题 11

除性分式理想：VII.1.1

等价除子：VII.1.2

整环，Bézout（或 Bezout）：VII.1.习题 20

--- 完全整闭整环：V.1.4

--- Dedekind 整环：VII.2.1

--- 因子分解整环：VII.3.1

--- 整闭整环：V.1.2

--- 有限特征的整闭整环：VII.1.习题 25、26 和 28

--- 整 Noether 整环：V.3.习题 6

--- Krull 整环：VII.1.3

--- 一维局部整环：VI.4.习题 7

--- Prüfer 整环（或 Prüfer）：VII.2.习题 12

--- 伪 Bézout 整环：VII.1.习题 21

--- 伪主理想整环：VII.1.习题 21

--- 伪 Prüfer 整环：VII.2.习题 19

--- 正则整闭整环：VII.1.习题 30

支配（局部环）局部环：VI.1.1

模的代数托里克对偶：VI.5.习题 9

--- 格的对偶：VII.4.2

--- 拓扑托里克对偶：VI.5.习题 10

拓扑幂零元：III.4.3

代数相关、无关的元素：III.1.1

--- 强互素的元素：III.4.1

殆幂零自同态：IV.1.4

等价除子：VII.1.2

--- 等价赋值：VI.3.2

本性分次理想：III.1.4

--- 本性赋值：VII.1.4

欧氏有序域：VI.2.习题 4

穷竭滤过：III.2.1

拟 Galois 扩张：V.2.2

不变因子：VII.4.习题 11 和 14

因子分解整环：VII.3.1

赋值的典范分解：VI.3.2

忠实平坦模：I.3.1

代数自由、相关的族：III.1.1

--- 形式自由的：III.2.9

在代数中代数闭的域：V.1.2

--- 分解域：V.2.3

--- 射影域：VI.2.1

--- 局部环的剩余域：II.3.1

位的剩余域：VI.2.3

--- 赋值的剩余域：VI.3.2

--- 位的值域：VI.2.2

滤过群、环、模：III.2.1

m-进滤过：III.2.1

--- 与分次相伴的：III.2.1

--- 与环结构、模结构相容的：III.2.1

--- 离散的：III.2.5

--- m-好的：III.3.1

--- 递增、递减、分离、穷竭滤过：III.2.1

--- 诱导的、乘积的、商的：III.2.1

--- 由环滤过导出的模滤过：III.2.1

--- 平凡的：III.2

环上的有限代数：V.1.1

--- （位）在元素处有限：VI.2.2

有限生成代数：III.1.1

有限表示的：I.2.8

对 M 平坦、M-平坦（模）：I.2.2

--- 平坦模：I.2.3

形式自由族：III.2.9

分式理想：VII.1.1

阶函数：III.2.2

Gauss 整数：V.1.1

Gauss 引理：VII.3.5

Gelfand-Mazur 定理：VI.6.4.

由子集生成（乘性子集）：II.2.1

生成元的形式系：III.2.9

m-好滤过：III.3.1

分解群：V.2.2

--- 滤过群：III.2.1

--- 惯性群：V.2.2

--- 可逆模类群：II.5.7

--- 局部有限算子群：V.1.9

--- 作用在环上的群：V.1.9

--- 赋值的阶群：VI.3.2

--- 高度 n、高度 +∞ 的有序群：VI.4.4

高度 \(\leqslant 1\)（素理想的）：VII.1.6

--- 有序群、赋值的高：VI.4.4

Hensel 环：III.4.习题 3

Hensel 条件：III.4.5

--- 定理：III.4.3

素理想 \(p'\) 的分解群的典范同态

--- 从 A' 到 A'/p' 的自同构群：V.2.2

--- 与滤过相容：III.2.4

--- 与滤过相容的同态相伴的分次同态：III.2.4

--- 局部同态：II.3.1

--- 伪单射、伪满射、伪零、伪双射：VII.4.4

行列式理想：VII.4.习题 10 和 14

--- 本性分次理想：III.1.4

--- 浸入素理想：IV.2.3

--- 整理想、分式理想：VII.1.1

--- 可逆分式理想：II.5.7

--- 位于理想之上的理想：V.2.1

--- 极小素理想：II.2.6

--- 位的理想：VI.2.3

--- 赋值的理想：VI.3.2

--- 素理想：II.1.1

--- 与模相伴的素理想：IV.1.1

--- 完全分解的素理想：V.2.2

--- 高度 ≤1 的素理想：VII.1.6

--- 准素、p-准素理想：IV.2.1 和习题 20

--- 非分歧理想：V.2.习题 18 和 19

理想 Hausdorff 模：III.5.1

互素理想：II.1.2

Cauchy 恒等式：VII.3.习题 18

浸入素理想：IV.2.3

非真赋值：VI.3.1

递增滤过：III.2.1

独立赋值环：VI.7.2

--- 独立赋值：VI.7.2

有序群子群的初始指数、赋值的初始分歧指数：VI.8.4

--- 分歧指数：VI.8.1

诱导滤过：III.2.1

Noether 归纳法（原则）：II.4.2

惯性域：V.2.3

--- 惯性环、惯性群：V.2.2

初始分歧指数：VI.8.4

代数整数：V.1.1

--- Gauss 整数：V.1.1

--- 环上的整元：V.1.1

\begin{Shaded}
\begin{Highlighting}[]
\NormalTok{整数，n{-}进：III.2.12}
\NormalTok{环上的整代数：V.1.1}
\NormalTok{— 整闭包：V.1.2}
\NormalTok{— 整理想：VII.1.1}
\NormalTok{整闭整环：V.1.2}
\NormalTok{— 在代数（环）中整闭：V.1.2}
\NormalTok{一个格相对于另一个格的相对不变量：VII.4.6}
\NormalTok{可逆分式理想：II.5.7}
\NormalTok{— 可逆子模：II.5.6}
\NormalTok{不可约分支：II.4.1}
\NormalTok{— 不可约集：II.4.1}
\NormalTok{— 不可约空间：II.4.1}
\NormalTok{孤立子群：VI.4.2}
\NormalTok{Jacobson 环：V.3.4}
\NormalTok{Krull{-}Akizuki 定理：VII.3.5}
\NormalTok{Krull 整环：VII.1.3}
\NormalTok{Krull\textquotesingle{} 定理：III.3.2}
\NormalTok{Lasker 模：IV.2.习题 23}
\NormalTok{格：VII.4.1}
\NormalTok{— 对偶格：VII.4.2}
\NormalTok{— 自反格：VII.4.2}
\NormalTok{引理，Artin{-}Rees：III.3.1}
\NormalTok{— Gauss\textquotesingle{} 引理：VII.3.5}
\NormalTok{— 正规化引理：V.3.1}
\NormalTok{线性拓扑环：III.4.2}
\NormalTok{局部同态：II.3.1}
\NormalTok{— 局部环：II.3.1}
\NormalTok{局部有限算子群：V.1.9}
\NormalTok{位于理想之上的理想：V.2.1}
\NormalTok{优势集：VI.3.5}
\NormalTok{与环同态相伴的连续映射：II.4.3}
\NormalTok{极小多项式：V.1.3}
\NormalTok{— 极小素理想：II.2.6}
\NormalTok{凝聚模：I.2.习题 11}
\NormalTok{— 忠实平坦模：I.3.1}
\NormalTok{— 滤过模：III.2.1}
\NormalTok{— 与分次模相伴的滤过模：III.2.1}
\NormalTok{— 有限表示模：I.2.8}
\NormalTok{— 平坦模：I.2.3}
\NormalTok{— 对 M 平坦、M{-}平坦：I.2.2}
\end{Highlighting}
\end{Shaded}

\begin{Shaded}
\begin{Highlighting}[]
\NormalTok{— 极小：V.1.3}
\end{Highlighting}
\end{Shaded}

与滤过模相伴的分次模：III.2.3

--- 理想 Hausdorff 模：III.5.1

--- 由环的子集定义的分式模：II.2.2

--- 秩 n 投射模：II.5.3

--- 伪凝聚模：I.2.习题 11

--- 伪零模：VII.4.4

除子类幺半群：VII.1.2

非处处有定义的合成律的态射：VI.2.1

乘性子集：II.2.1

环的幂零根：II.2.6

Noether 空间：II.4.2

非退化子模：II.5.5

正规化引理：V.3.1

赋范离散赋值：VI.3.6

零点定理（Nullstellensatz）：V.3.3

赋值的阶群：VI.3.2

--- 元素在赋值下的阶：VI.3.2

--- 形式幂级数的约化阶：VII.3.8

具线性扩张性质的环对：I.3.7

Ostrowski 定理：VI.6.4

位，在 \(x\) 处有限：VI.2.2

--- 平凡位：VI.2.2

不可约空间的泛点：II.4.习题 2

Newton 多边形：VI.4.习题 11

特异多项式：VII.3.8

预备定理：VII.3.8

模的表示，--- 有限表示：I.2.8

\(n\)-表示：I.2.习题 6

有限表示的（模）：I.2.8

准素分解：IV.2.2 和习题 20

--- p-准素理想、p-准素子模：IV.2.1 和习题 20

素理想：II.1.1

--- 素谱：II.4.3

主除子：VII.1.1

--- 主限制阿代尔：VII.2.4

Noether 归纳原理：II.4.2

乘积滤过：III.2.1

射影域：VI.2.1

伪单射、伪满射、伪零、伪双射（同态）：VII.4.4

伪同构：VII.4.4

伪零模：VII.4.4

拟 Galois 扩张：V.2.2

商滤过：III.2.1

理想的根：II.2.6 投射模在 p 处的秩：II.5.3 --- 投射模的秩：II.5.3 --- 交换群的有理秩：VI.10.2 --- 剩余秩：VI.8.5 交换群的有理秩：VI.10.2 约化阶：VII.3.8 --- 准素分解：IV.2.3 --- 约化环：II.2.6 --- 约化幂级数：VII.3.8 自反格：VII.4.2 相关联的局部环：VI.1.习题 1 互素理想：II.1.2 极端元的代表系：VII.3.3 赋值的剩余类次数：VI.8.1 --- 剩余域：II.3.1 --- 赋值的剩余秩：VI.8.5 模的有限自由分解：VII.4.7 限制阿代尔：VII.2.4 --- 限制形式幂级数：III.4.2 绝对平坦环：I.2.习题 17 --- （左、右）凝聚环：I.2.习题 17 --- 分解环：V.2.2 --- 滤过环：III.2.1 --- 与分次环相伴的滤过环：III.2.1 --- 与滤过环相伴的分次环：III.2.3 --- 惯性环：V.2.2 --- 在代数中整闭：V.1.2 --- Jacobson 环：V.3.4 --- 线性拓扑环：III.4.2 --- 局部环：II.3.1 --- 支配局部环的局部环：VI.1.1 --- A 在 p 处的局部环、p 的局部环（p 为素理想）：II.3.1 --- 位的环：VI.2.3 --- 赋值的环：VI.3.2

由环的子集定义的分式环：II.2.1

--- 约化环：II.2.6

--- 半局部环：II.3.5

--- 全分式环：II.2.1

--- 非分歧环：V.2.习题 19

--- 赋值环，域的赋值环：VI.1.1

--- Zariski 环：III.3.3

独立赋值环：VI.7.2

子模关于乘性子集的饱和（关于

素理想）：II.2.4

半局部环：II.3.5

约化幂级数：VII.3.8

--- 限制形式幂级数：III.4.2

不可约集：II.4.1

--- 全序群中的优势集：VI.3.5

不可约空间：II.4.1

--- Noether 空间：II.4.2

特殊拓扑：II.4.3

环的素谱：II.4.3

强 Lasker 模：IV.2.习题 28

--- 强准素子模：IV.2.习题 27

--- 强互素的元素：III.4.1

有序群的孤立子群：VI.4.2

可逆子模：II.5.6

--- 非退化子模：II.5.5

--- 准素、p-准素子模：IV.2.1

环的乘性子集：II.2.1

--- 由子集生成的乘性子集：II.2.1

--- 饱和乘性子集：II.2.习题 1

模的支撑集：II.4.4

赋值扩张的完备系：VI.8.2

--- 生成元的形式系：III.2.9

--- 极端元的代表系：VII.3.3

绝对值逼近定理：VI.7.3

--- 赋值逼近定理：VI.7.2

--- Gelfand-Mazur 定理：VI.6.4

--- Hensel 定理：III.4.3

--- Hilbert 零点定理（Nullstellensatz）：V.3.3

--- Krull 定理：III.3.1

--- Krull-Akizuki 定理：VII.3.5

--- Ostrowski 定理：VI.6.3

预备定理：VII.3.8

--- Stickelberger 定理：VI.8.习题 18

--- Zariski 主定理：V.3.习题 7

拓扑幂零元：III.4.3

由滤过定义的拓扑：III.2.5

--- 谱拓扑：II.4.3

--- Zariski 拓扑：II.4.3

传递子：I.2.10

平凡滤过：III.2.1

--- 平凡位：VI.2.2

超度量绝对值：VI.6.1

离散赋值的一致化元：VI.3.6

非分歧赋值：VI.8.1

赋值、元素 x 的赋值：VI.3.1 和 2

--- 离散、赋范离散赋值：VI.3.6

--- 本性赋值：VII.1.4 和习题 26

--- 非真赋值：VI.3.1

--- 赋值环：VI.1.1

--- 非分歧赋值：VI.8.1

等价赋值：VI.3.2

--- 独立赋值：VI.7.2

位的值域：VI.2.2

--- 超度量绝对值：VI.6.1

与模相伴的弱伴随素理想：IV.1.习题 17

Zariski 环：III.3.3

--- Zariski 拓扑：II.4.3

